\chapter{Epílogo: La Semilla de 2102}

\neofigure{12_epilogo_semilla_2102.jpg}{Reunión secreta del Directorio de Orekak-Corp analizando la plasticidad cerebral infantil.}

\section{El amanecer sobre la Ría de Bilbao}

La sobrecarga del búnker subterráneo se disipa en un estallido sordo de luz blanca que apaga los neones del Distrito 7 durante varios minutos. Los personajes emergen empapados en las dársenas exteriores mientras el humo negro se confunde con el sirimiri matutino.

\begin{itemize}
  \item \lead{El destino de Ane Garate:} Aunque el puerto neuronal ha quedado inerte, las cicatrices psicológicas son profundas. Con el cuidado de su familia y la atención clandestina de Maite «Sorgin» y el Celador, Ane recupera lentamente la memoria y la voz propia, aunque jamás volverá a mirar las pantallas de televisión corporativas sin estremecerse.
  \item \lead{La reacción en La Margen:} Si los datos fueron transmitidos, las tascas del Distrito 7 hierven de indignación y huelgas salvajes. Orekak-Corp despliega cordones policiales y anuncia en comunicados urgentes que el Proyecto IZAR fue «un sabotaje terrorista de elementos marginales contra el legítimo Programa de Excelencia». El programa es cancelado oficialmente de la noche a la mañana.
\end{itemize}

\section{En la penumbra del Directorio (Solo para el Director)}

Tres días después del colapso del Proyecto IZAR, en una sala acorazada en lo alto de la Cúpula Dorada, tres altos directivos de Orekak-Corp revisan los registros médicos sobre una mesa de cristal negro:

\begin{quote}
\textit{«El informe de autopsia es inequívoco, señores directores. El cerebro adulto es un callejón sin salida. Posee demasiada biografía, demasiados apegos, demasiada resistencia biológica al mandato colectivo. No toleran el bus de datos.}

\textit{Si queremos que el sistema operativo de NeoBilbao funcione, necesitamos un sustrato biológico virgen. Mentes que no hayan conocido el mundo exterior. Mentes con una neuroplasticidad maleable, capaces de aceptar una simulación de paraíso sin cuestionar el cable que llevan en la nuca.}

\textit{Para la próxima generación... no reclutaremos adultos. Nos quedaremos con sus hijos recién nacidos.»}
\flushright --- Acta reservada del Directorio de Orekak-Corp (Diciembre de 2084).
\end{quote}

\section{De 2084 a 2102: El puente entre campañas}

Para los grupos que deseen enlazar este módulo con la campaña principal de \textbf{NeoBilbao: Etxearen Zina (2102)}:

\begin{itemize}
  \item \lead{Los ecos del pasado:} Los personajes que sobrevivieron a 2084 son ahora los ancianos amargados que en 2102 advierten en susurros a los jóvenes mineros que no se fíen del Velo.
  \item \lead{El linaje de la resistencia:} Los hijos o sobrinos de los protagonistas de 2084 serán los obreros que dieciocho años después trabajarán en la fundición de Burdin-Urdin, creyendo ciegamente en el Juramento del Linaje hasta que el cortocircuito del primer capítulo de 2102 les obligue a despertar.
  \item \lead{La continuidad temática:} En 2084 lucharon por la verdad en las calles de asfalto y neón; en 2102 sus descendientes lucharán por sus hijos en el abismo de hierro de los Zócalos Industriales.
\end{itemize}
